\documentclass[11pt]{article}
\usepackage[margin=0.88in]{geometry}
\usepackage{amsmath,amssymb,amsthm,graphicx,microtype,xcolor}
\usepackage[hidelinks]{hyperref}
\setlength{\parindent}{0pt}
\setlength{\parskip}{0.52em}
\newtheorem{theorem}{Theorem}
\newtheorem{lemma}[theorem]{Lemma}
\newtheorem{corollary}[theorem]{Corollary}
\newcommand{\F}{\mathbb F}
\newcommand{\wt}{\operatorname{wt}}
\newcommand{\supp}{\operatorname{supp}}
\newcommand{\status}{\textcolor{blue!60!black}{\textbf{candidate full result---likely valid}}}

\begin{document}

\begin{center}
{\Large\bfseries An infinite-coordinate MacWilliams theorem and general
convolutional-code isometries}\par
\vspace{0.35em}
{\large Full-result packet for the open problem in arXiv:1502.02635}\par
\vspace{0.45em}
\status
\end{center}

\section*{1. Source problem}

Ferrer, Gary, and Hern\'andez ask for a representation of general
$\F$-isometries between convolutional codes in the introduction of
\cite{FGH}; the statement occurs on page~2 of arXiv:1502.02635v1:

\begin{center}
\includegraphics[width=\textwidth]{figures/open_problem_crop.png}
\end{center}

In exact terms, the target sentence is:
\begin{quote}
``It remains open the representation of general $\F$-isometries defined
between convolutional codes.''
\end{quote}
Here ``general'' contrasts with the $\F[z]$-module isometries treated in
\cite{GL}: the map is only required to be $\F$-linear and to preserve the
coefficientwise Hamming weight.  If
$v(z)=\sum_{t\geq0}v^{(t)}z^t\in\F[z]^n$, that weight is
\[
  \wt(v)=\sum_{t\geq0}\wt_{\F^n}(v^{(t)}).
\]

We prove a more general extension theorem for arbitrary subspaces of an
algebraic direct sum.  Applied to $\F[z]^n$, it gives the requested complete
representation: a general isometry is precisely a restriction of a weighted
permutation of the active time--coordinate pairs.  No controllability,
$\F[z]$-linearity, degree preservation, or finite-dimensionality is assumed.

\section*{2. Definitions and proof intuition}

For a set $I$, write
\[
 \F^{(I)}=\{x:I\to\F: \supp(x)\text{ is finite}\},
 \qquad \wt(x)=|\supp(x)|.
\]
If $V\leq\F^{(I)}$, its \emph{active coordinates} are
\[
 I_V=\{i\in I:\text{some }x\in V\text{ has }x(i)\ne0\}.
\]
Two nonzero linear functionals are \emph{projectively equivalent} if one is a
nonzero scalar multiple of the other.

\textbf{Proof intuition.}
Every coordinate of a code restricts to a linear functional on the abstract
code space.  Weight is exactly the number of these functionals that do not
vanish at a given vector.  On a finite-dimensional subspace, only finitely
many coordinate functionals remain nonzero, and a finite incidence-matrix
argument says that the weight function determines their projective
multiplicities.  The apparent infinite-dimensional gluing problem disappears
because the coordinate families are pointwise finite: one vector sees only
finitely many functionals.  A finite-dimensional subspace can therefore
isolate any chosen global projective class.  The source and target coordinate
families must match class by class, which is exactly a weighted coordinate
permutation.

\section*{3. A finite uniqueness lemma}

\begin{lemma}[projective functional multiplicities]\label{lem:finite}
Let $E$ be a finite-dimensional vector space over $\F_q$.  Let
$\mathcal L$ and $\mathcal M$ be finite multisets of nonzero members of $E^*$.
If
\[
 \#\{\ell\in\mathcal L:\ell(x)\ne0\}
 =\#\{m\in\mathcal M:m(x)\ne0\}
 \qquad(x\in E),
\]
then $\mathcal L$ and $\mathcal M$ have the same multiplicity in every
projective class of $E^*$.
\end{lemma}

\begin{proof}
The assertion is immediate when $\dim E\leq1$, so suppose $d=\dim E\geq2$.
Index the rows of a matrix $A$ by the projective points $P=[x]$ of $E$ and the
columns by the projective nonzero functionals $H=[\ell]$ in $E^*$, and put
\[
 A_{P,H}=\mathbf 1_{\{\ell(x)\ne0\}}.
\]
There are $(q^d-1)/(q-1)$ rows and the same number of columns.  If $P=Q$,
there are $q^{d-1}$ projective functionals nonzero on $P$.  If $P\ne Q$,
inclusion--exclusion among the hyperplanes through $P$ and $Q$ gives
\[
 (AA^{\mathsf T})_{P,Q}=q^{d-2}(q-1).
\]
Thus, with $J$ the all-ones matrix,
\[
 AA^{\mathsf T}
 =q^{d-2}I+q^{d-2}(q-1)J.
\]
This is positive definite over $\mathbb R$, so $A$ is invertible.  If $c_H$
is the difference between the two multiplicities in class $H$, the assumed
weight equality says $Ac=0$.  Hence $c=0$.
\end{proof}

This is the finite MacWilliams uniqueness mechanism in the exact form needed
below; its proof has been included so that the infinite-dimensional theorem
does not depend on an extension theorem with an implicit finite-length
hypothesis.

\section*{4. Infinite-coordinate extension theorem}

\begin{theorem}[finitary MacWilliams extension]\label{thm:main}
Let $I,J$ be arbitrary sets, let $V\leq\F_q^{(I)}$ and
$W\leq\F_q^{(J)}$, and let $T:V\to W$ be an $\F_q$-linear bijection.  Then
the following are equivalent.
\begin{enumerate}
\item $\wt(Tx)=\wt(x)$ for every $x\in V$.
\item There are a bijection $\sigma:J_W\to I_V$ and scalars
      $a_j\in\F_q^*$ such that
      \begin{equation}\label{eq:monomial}
        (Tx)(j)=a_jx(\sigma(j))
        \qquad(x\in V,\ j\in J_W).
      \end{equation}
      On $J\setminus J_W$, every member of $W$ is zero by definition.
\end{enumerate}
In particular, after puncturing inactive coordinates, every linear Hamming
isometry between possibly infinite-dimensional finitary codes is the
restriction of a monomial coordinate transformation.
\end{theorem}

\begin{proof}
For $i\in I_V$, let $\alpha_i\in V^*$ be evaluation at $i$.  For
$j\in J_W$, pull target evaluation back to $V$:
\[
 \beta_j(x)=(Tx)(j).
\]
Both families are \emph{pointwise finite}: for fixed $x\in V$, only finitely
many $\alpha_i(x)$ are nonzero, and only finitely many $\beta_j(x)$ are
nonzero because $Tx\in\F_q^{(J)}$.  Weight preservation is precisely
\begin{equation}\label{eq:weights}
 \#\{i:\alpha_i(x)\ne0\}=\#\{j:\beta_j(x)\ne0\}
 \qquad(x\in V).
\end{equation}

First restrict to a finite-dimensional subspace $E\leq V$.  Choose a finite
basis of $E$.  Pointwise finiteness shows that only finitely many
$\alpha_i|_E$ are nonzero: their indices lie in the union of the finitely many
supports of the basis vectors.  The same holds for the $\beta_j|_E$.
Equation~\eqref{eq:weights} and Lemma~\ref{lem:finite} show that the two
restricted families have equal projective multiplicities.

We now recover the global classes from these local statements.  Fix a
nonzero functional $\alpha$ occurring in the first family and choose
$x_0\in V$ with $\alpha(x_0)\ne0$.  Only finitely many members $\varphi$ of
the union of the two families satisfy $\varphi(x_0)\ne0$.  For every such
$\varphi$ not globally proportional to $\alpha$, choose $y_\varphi\in V$
for which
\[
 \alpha(x_0)\varphi(y_\varphi)
 \ne \varphi(x_0)\alpha(y_\varphi).
\]
Let $E$ be the span of $x_0$ and these finitely many $y_\varphi$.  A family
member whose restriction to $E$ is projectively equivalent to $\alpha|_E$
must be nonzero at $x_0$, and the chosen separator then shows that it is
globally proportional to $\alpha$.  Conversely every global multiple plainly
has that restricted class.  The local multiplicity equality on $E$ therefore
says that the global multiplicity of $[\alpha]$ is the same in the two
families.

This treats every class occurring among the $\alpha_i$.  Repeating the same
argument with the two families interchanged rules out a class occurring only
among the $\beta_j$.  Hence the multisets
\[
 \{[\alpha_i]:i\in I_V\}
 \quad\text{and}\quad
 \{[\beta_j]:j\in J_W\}
\]
are equal, with multiplicity.  Match equal classes to obtain a bijection
$\sigma:J_W\to I_V$ and choose $a_j\ne0$ with
$\beta_j=a_j\alpha_{\sigma(j)}$.  This is exactly
\eqref{eq:monomial}.  The converse follows immediately by counting nonzero
coordinates.
\end{proof}

\section*{5. Full classification for polynomial convolutional codes}

\begin{corollary}[general $\F$-isometries of convolutional codes]
\label{cor:conv}
Let $\mathcal C\leq\F_q[z]^n$ and
$\mathcal D\leq\F_q[z]^m$ be polynomial convolutional codes, and let
$\Phi:\mathcal C\to\mathcal D$ be an arbitrary $\F_q$-linear bijection.
For a code $\mathcal C$, define its active time--coordinate support by
\[
 S_{\mathcal C}=\{(t,r)\in\mathbb N_0\times\{1,\dots,n\}:
 [z^t]c_r\ne0\text{ for some }c\in\mathcal C\}.
\]
Then $\Phi$ preserves Hamming weight if and only if there are a bijection
$\tau:S_{\mathcal D}\to S_{\mathcal C}$ and nonzero scalars
$u_{t,r}\in\F_q^*$ such that, writing $\tau(t,r)=(s,k)$,
\begin{equation}\label{eq:conv}
 [z^t](\Phi c)_r=u_{t,r}[z^s]c_k
 \qquad(c\in\mathcal C,\ (t,r)\in S_{\mathcal D}).
\end{equation}
All coefficients outside $S_{\mathcal D}$ vanish.
\end{corollary}

\begin{proof}
Identify $\F_q[z]^n$ with
$\F_q^{(\mathbb N_0\times\{1,\dots,n\})}$ by taking coefficients.  The
convolutional weight is exactly the resulting finitary Hamming weight, so
Theorem~\ref{thm:main} applies and gives~\eqref{eq:conv}.  The converse is
again immediate.
\end{proof}

Corollary~\ref{cor:conv} is a representation theorem for the general
$\F$-isometries left open in \cite{FGH,GL}.  It also explains why a general
map need not resemble an $\F[z]$-module map: the bijection $\tau$ may depend
arbitrarily on time and may mix time and component coordinates.  Module
linearity is the additional condition that this coordinate rule intertwines
multiplication by $z$.

\section*{6. Verification, scope, and novelty boundary}

The incidence counts in Lemma~\ref{lem:finite} were independently enumerated
for prime fields $q\in\{2,3,5\}$ and dimensions $1\leq d\leq4$ by
\texttt{code/verify\_incidence.py}.  The script checks all projective point
pairs, the stated diagonal and off-diagonal Gram entries, and the resulting
positive eigenvalues.  This is a sanity check, not a substitute for the
symbolic proof.

The two proof-sensitive points for human review are: (i) pointwise finiteness
really does imply finiteness after restriction to a finite-dimensional
subspace, by taking the union of the supports of a basis; and (ii) the finite
separator space isolates a \emph{global} projective class simultaneously in
both families.  Both steps are explicit in the proof.

The result fully answers the quoted polynomial convolutional-code
representation problem.  It does \emph{not} answer Question~1.4 of \cite{FGH}
for arbitrary nondiscrete measured locally compact spaces, and it does not
classify isometries of Laurent-series codes whose words may have infinite
support or infinite weight.

A bounded novelty search on 11 August 2026 covered the run's local indexes,
the exact source title and quoted open sentence, ``general $\F$-isometries of
convolutional codes'', ``infinite-length MacWilliams extension'', ``finitely
supported Hamming isometry'', and nearby variants.  It also checked
arXiv:0902.2235 (module isometries), arXiv:1507.05212 (finite-length geometric
extension results), and arXiv:1709.06070 (infinite alphabet rings rather than
infinite coordinate length).  No exact infinite-coordinate theorem or later
answer to the quoted problem was found.  Novelty confidence is therefore
moderate, not definitive: the short argument could exist as folklore or under
different terminology.  Expert literature review is recommended before a
publication-level novelty claim.

\small
\begin{thebibliography}{9}
\bibitem{FGH}
M. Ferrer, M. Gary, and S. Hern\'andez,
\emph{Weight-preserving isomorphisms between spaces of continuous functions:
the scalar case}, arXiv:1502.02635v1 (2015); J. Math. Anal. Appl. 433 (2016),
223--235.

\bibitem{GL}
H. Gluesing-Luerssen,
\emph{On isometries for convolutional codes}, Adv. Math. Commun. 3 (2009),
179--203; arXiv:0902.2235.

\bibitem{MacWilliams}
F. J. MacWilliams,
\emph{A theorem on the distribution of weights in a systematic code},
Bell System Tech. J. 42 (1963), 79--94.
\end{thebibliography}

\end{document}
